\chapter{Cálculo diferencial}\label{ch:b2:diffcalc}

O cálculo a duas variáveis do volume do primeiro ano de graduação amadurece no
cálculo diferencial de aplicações entre espaços normados: a
\emph{\omterm{def:b2:diffcalc:differential}{diferencial}} como melhor aproximação linear, a regra da cadeia
em plena generalidade, o teorema da simetria de Schwarz \emph{demonstrado}, fórmulas
de Taylor e a análise completa de segunda ordem dos extremos. O
teorema da função inversa, coroa da teoria, é enunciado com sua
estratégia de demonstração --- um ponto fixo de Banach --- explicitada.

Em todo o capítulo, $U$ é um \omterm{def:b2:metric:topology}{aberto} de $\R^n$ (ou de um espaço normado; o
caso de dimensão finita carrega todas as ideias), $f \colon U \to
\R^m$.

\section{A diferencial}

\begin{definition}\label{def:b2:diffcalc:differential}
$f$ é \emph{diferenciável}\index{diferenciável} em $a$ quando existe
uma aplicação linear (\omterm{def:b2:metric:continuity}{contínua}) $\dd f_a \colon \R^n \to \R^m$ com
\[
f(a + h) = f(a) + \dd f_a(h) + o(\norm h)
\qquad (h \to 0).
\]
A aplicação $\dd f_a$, a \emph{diferencial}\index{diferencial} de
$f$ em $a$, é única; sua matriz nas bases canônicas é
a \emph{matriz jacobiana}\index{matriz jacobiana} $J_f(a) =
\bigl(\frac{\partial f_i}{\partial x_j}(a)\bigr)$.
A diferenciabilidade implica a \omterm{def:b2:metric:continuity}{continuidade} e a existência de todas as
derivadas direcionais $\dd f_a(v) = \lim_{t\to0}\frac{f(a + tv)
- f(a)}{t}$; a recíproca falha
(\cref{exo:b2:diffcalc:2}). Para $m = 1$, $\dd f_a(h) = \langle
\nabla f(a), h\rangle$: o gradiente do primeiro ano, agora entendido como
o vetor que representa a diferencial.
\end{definition}

\begin{theorem}[Critério $C^1$]\label{thm:b2:diffcalc:c1}
Se todas as derivadas parciais de $f$ existem em $U$ e são \omterm{def:b2:metric:continuity}{contínuas}
em $a$, então $f$ é \omterm{def:b2:diffcalc:differential}{diferenciável} em $a$. ``$C^1$ em $U$'' ---
parciais \omterm{def:b2:metric:continuity}{contínuas} --- implica portanto a \omterm{def:b2:diffcalc:differential}{diferenciabilidade}
em toda parte, com \omterm{def:b2:diffcalc:differential}{diferencial} \omterm{def:b2:metric:continuity}{contínua}.
\end{theorem}

\begin{proof}
Componente a componente ($m = 1$ basta). A demonstração do primeiro ano para duas
variáveis --- mova uma coordenada de cada vez, aplique o teorema do valor
médio de uma variável em cada trecho e use a \omterm{def:b2:metric:continuity}{continuidade} das parciais
em $a$ --- generaliza-se palavra por palavra a $n$ trechos:
\[
f(a + h) - f(a) = \sum_{j=1}^{n} \bigl(f(a + h^{(j)}) - f(a +
h^{(j-1)})\bigr)
= \sum_j h_j\,\frac{\partial f}{\partial x_j}(\xi_j),
\]
em que $h^{(j)}$ congela as $j$ primeiras coordenadas de $h$ e
$\xi_j$ está no $j$-ésimo trecho; a \omterm{def:b2:metric:continuity}{continuidade} transforma cada
$\frac{\partial f}{\partial x_j}(\xi_j)$ em $\frac{\partial
f}{\partial x_j}(a) + o(1)$, e o erro é $o(\norm h)$.
\end{proof}

\begin{theorem}[Regra da cadeia]\label{thm:b2:diffcalc:chainrule}
Se $f$ é \omterm{def:b2:diffcalc:differential}{diferenciável} em $a$ e $g$ em $f(a)$, então $g \circ f$
é \omterm{def:b2:diffcalc:differential}{diferenciável} em $a$ com
\[
\dd(g \circ f)_a = \dd g_{f(a)} \circ \dd f_a ,
\qquad
J_{g\circ f}(a) = J_g\bigl(f(a)\bigr)\,J_f(a) :
\]
as jacobianas se multiplicam.\index{regra da cadeia}
\end{theorem}

\begin{proof}
Escreva $f(a + h) = f(a) + \dd f_a(h) + \norm h\,\varepsilon_1(h)$
e $g(b + k) = g(b) + \dd g_b(k) + \norm k\,\varepsilon_2(k)$ com
$b = f(a)$, $k = k(h) = \dd f_a(h) + \norm h \varepsilon_1(h)$.
Substituindo,
\[
g(f(a+h)) = g(b) + \dd g_b\bigl(\dd f_a(h)\bigr)
+ \norm h\,\dd g_b(\varepsilon_1(h)) + \norm{k}\,\varepsilon_2(k),
\]
e os dois termos de erro são $o(\norm h)$: o primeiro porque $\dd g_b$ é
\omterm{def:b2:metric:continuity}{contínua} e $\varepsilon_1 \to 0$; o segundo porque $\norm k
\leq C\norm h$ (aplicação linear limitada mais termo pequeno) e
$\varepsilon_2(k) \to 0$ quando $h \to 0$.
\end{proof}

\begin{example}[Funções radiais, de uma vez por todas]\label{ex:b2:diffcalc:radial}
Sejam $r(x) = \norm x_2$ em $\R^n\setminus\{0\}$ e $f = g \circ
r$ com $g$ uma função $C^1$ de uma variável. Primeiro, $r$ é
\omterm{def:b2:diffcalc:differential}{diferenciável} fora de $0$: de $r^2 = \sum x_i^2$,
\[
\frac{\partial r}{\partial x_i} = \frac{x_i}{r},
\qquad\text{i.e.}\qquad
\nabla r(x) = \frac{x}{\norm x} ,
\]
o vetor radial \omterm{def:b2:hermitian:adjoint}{unitário} (derive $r^2$ e divida --- ou
aplique a regra da cadeia a $\sqrt{\cdot}$). Então a regra da cadeia
dá, para toda função radial,
\[
\nabla f(x) = g'\bigl(\norm x\bigr)\,\frac{x}{\norm x} .
\]
Instância trabalhada: $g(r) = \frac1r$ fornece $\nabla\frac{1}{\norm
x} = -\frac{x}{\norm x^3}$, o campo inverso do quadrado da
gravitação e da eletrostática --- direção radial, magnitude
$\frac{1}{\norm x^2}$. Lição final: os gradientes de funções radiais
são radiais porque as superfícies de nível são esferas e
o gradiente é ortogonal às superfícies de nível; em $x = 0$, em
contraste, $r$ \emph{não} é \omterm{def:b2:diffcalc:differential}{diferenciável} (nenhuma aplicação
linear candidata casa com $\norm h$ vinda de todas as direções) --- perfis
radiais suaves precisam de $g'(0) = 0$ para cruzar a origem
com elegância.
\end{example}

\begin{theorem}[Desigualdade do valor médio]\label{thm:b2:diffcalc:mvi}
Seja $f$ \omterm{def:b2:diffcalc:differential}{diferenciável} em $U$ e suponha que o segmento $\intcc{a}{b}
= \{a + t(b-a)\}$ esteja em $U$. Então
\[
\norm{f(b) - f(a)} \leq \norm{b - a}\,
\sup_{x \in \intcc{a}{b}} \vertiii{\dd f_x} .
\]
Em particular, uma aplicação \omterm{def:b2:diffcalc:differential}{diferenciável} de \omterm{def:b2:diffcalc:differential}{diferencial} nula num
\omterm{def:b2:metric:topology}{aberto} \emph{\omterm{def:b2:metric:connected}{conexo}} é constante.
\end{theorem}

\begin{proof}
A função $\varphi(t) = f(a + t(b-a))$ é \omterm{def:b2:diffcalc:differential}{diferenciável} em
$\intcc{0}{1}$ com $\varphi'(t) = \dd f_{a + t(b-a)}(b - a)$
(regra da cadeia), de \omterm{def:b2:nvs:norm}{norma} $\leq M\norm{b-a}$ com $M$ o sup
exibido. Para $f$ com valores em $\R$, a desigualdade do valor médio de uma variável
conclui; para valores vetoriais, aplique-a a $t \mapsto \langle u,
\varphi(t)\rangle$ com $u$ o vetor unitário ao longo de $f(b) - f(a)$.
Constância: localmente constante (segmentos em bolas) mais \omterm{def:b2:metric:connected}{conexidade}
(o conjunto em que $f$ vale um dado valor é \omterm{def:b2:metric:topology}{aberto} e fechado:
\cref{ch:b2:metric}).
\end{proof}

\begin{example}[Uma constante de Lipschitz pela DVM]\label{ex:b2:diffcalc:lipschitz}
$f(x, y) = \sin x\,\sin y$ é \omterm{def:b2:metric:continuity}{lipschitziana} em $\R^2$, e com
qual constante? Seu gradiente é $\nabla f = (\cos x\sin y,\
\sin x\cos y)$, de \omterm{def:b2:nvs:norm}{norma} ao quadrado
\[
\cos^2x\sin^2y + \sin^2x\cos^2y
\leq \sin^2 y + \cos^2y\cdot 1 = 1
\]
(majore $\cos^2x$ e $\sin^2x$ por $1$ separadamente), logo
$\vertiii{\dd f_{(x,y)}} = \norm{\nabla f} \leq 1$ em toda parte,
e \cref{thm:b2:diffcalc:mvi} no segmento entre dois pontos
quaisquer dá
\[
\abs{f(b) - f(a)} \leq \norm{b - a}_2 :
\]
$f$ é $1$-lipschitziana, e a constante é ótima (perto da
origem, $f(x, \tfrac\pi2) = \sin x$ tem inclinação $1$). Lição
final: a desigualdade do valor médio converte uma cota \omterm{def:b2:funcseq:def}{pontual}
sobre a \omterm{def:b2:diffcalc:differential}{diferencial} num módulo global de \omterm{def:b2:metric:continuity}{continuidade} ---
a rota padrão para estimativas \omterm{def:b2:metric:continuity}{lipschitzianas} em qualquer dimensão,
e o motor dentro do~\cref{exo:b2:diffcalc:12}.
\end{example}

\section{Derivadas segundas}

\begin{theorem}[Schwarz]\label{thm:b2:diffcalc:schwarz}
Se $f$ é $C^2$ em $U$ (todas as parciais segundas existem e são
\omterm{def:b2:metric:continuity}{contínuas}), então, para todos $i, j$:\index{teorema de Schwarz}
\[
\frac{\partial^2 f}{\partial x_i\,\partial x_j}
= \frac{\partial^2 f}{\partial x_j\,\partial x_i} .
\]
\end{theorem}

\begin{proof}
Duas variáveis bastam ($x = x_i$, $y = x_j$, as outras congeladas).
Considere a diferença segunda
\[
\Delta(h) = f(a + h, b + h) - f(a + h, b) - f(a, b + h) + f(a,b) .
\]
Fixe $h$ e ponha $\varphi(x) = f(x, b+h) - f(x, b)$: então $\Delta(h)
= \varphi(a + h) - \varphi(a)$, e duas aplicações do teorema do valor
médio dão
\[
\Delta(h) = h\,\varphi'(\xi)
= h\Bigl(\frac{\partial f}{\partial x}(\xi, b+h) -
\frac{\partial f}{\partial x}(\xi, b)\Bigr)
= h^2\,\frac{\partial^2 f}{\partial y\,\partial x}(\xi, \eta),
\]
com $\xi \in \intoo{a}{a+h}$, $\eta \in \intoo{b}{b+h}$. Pela
\omterm{def:b2:metric:continuity}{continuidade}, $\frac{\Delta(h)}{h^2} \to \frac{\partial^2 f}{\partial
y\partial x}(a, b)$ quando $h \to 0$. O mesmo cálculo com os
papéis das variáveis trocados (congele primeiro a segunda variável)
dá $\frac{\Delta(h)}{h^2} \to \frac{\partial^2 f}{\partial x
\partial y}(a,b)$: os dois limites da mesma quantidade coincidem.
\end{proof}

\begin{example}[Por que $C^2$ é necessário: o contraexemplo de Peano]\label{ex:b2:diffcalc:peano}
Seja $f(x, y) = \dfrac{xy(x^2 - y^2)}{x^2 + y^2}$, $f(0,0) = 0$.
Fora da origem $f$ é $C^\infty$; na origem, todas as parciais primeiras
e segundas existem, mas as mistas discordam. Calcule
ao longo dos eixos: $f(x, 0) = f(0, y) = 0$ e, para $y \neq 0$,
\[
\frac{\partial f}{\partial x}(0, y)
= \lim_{x\to0}\frac{f(x,y)}{x}
= \frac{y(0 - y^2)}{y^2} = -y ,
\qquad\text{simetricamente}\qquad
\frac{\partial f}{\partial y}(x, 0) = x .
\]
Logo
\[
\frac{\partial^2 f}{\partial y\,\partial x}(0,0)
= \frac{\dd}{\dd y}\Bigl[\frac{\partial f}{\partial
x}(0,y)\Bigr]_{y=0} = -1,
\qquad
\frac{\partial^2 f}{\partial x\,\partial y}(0,0) = +1 :
\]
as duas parciais mistas existem e diferem. Nenhuma contradição com o
\cref{thm:b2:diffcalc:schwarz}: as parciais segundas de $f$ não são
\omterm{def:b2:metric:continuity}{contínuas} em $0$ (teste ao longo de $y = tx$). Lição final:
o teorema de Schwarz é um teorema genuíno sobre \emph{\omterm{def:b2:metric:continuity}{continuidade}},
e não uma identidade formal --- e a hipótese ``$C^2$'' no
Taylor--Young a seguir está fazendo trabalho de verdade.
\end{example}

\begin{theorem}[Taylor--Young de ordem 2]\label{thm:b2:diffcalc:taylor}
Seja $f \colon U \to \R$ de classe $C^2$ e $a \in U$. Então, quando $h \to 0$,
\[
f(a + h) = f(a) + \langle\nabla f(a), h\rangle
+ \frac12\, \langle H_a h,\, h\rangle + o\bigl(\norm h^2\bigr),
\]
em que $H_a = \bigl(\frac{\partial^2 f}{\partial x_i\partial
x_j}(a)\bigr)$ é a \emph{matriz hessiana}\index{matriz hessiana}
(\omterm{def:b2:quadratic:adjoint}{simétrica}, por Schwarz).
\end{theorem}

\begin{proof}
Aplique o teorema de Taylor--Young de uma variável (volume do primeiro ano de graduação) a
$\varphi(t) = f(a + th)$ em $\intcc{0}{1}$: pela regra da cadeia,
$\varphi'(t) = \langle \nabla f(a + th), h\rangle$ e $\varphi''(t)
= \langle H_{a+th}h, h\rangle$, ambas \omterm{def:b2:metric:continuity}{contínuas} em $t$. Então
$\varphi(1) = \varphi(0) + \varphi'(0) + \frac12\varphi''(\theta)$
(Taylor--Lagrange) com $\theta \in \intoo{0}{1}$, e a \omterm{def:b2:metric:continuity}{continuidade}
das parciais segundas converte $\varphi''(\theta) =
\varphi''(0) + o(1)\cdot\norm h^2$ \omterm{def:b2:funcseq:def}{uniformemente}: o
desenvolvimento exibido.
\end{proof}

\begin{example}[Um desenvolvimento de duas maneiras]\label{ex:b2:diffcalc:taylortwoways}
Desenvolva $f(x, y) = \eu^x\cos y$ na origem até a ordem $2$.
\emph{Por composição de desenvolvimentos de uma variável:}
\[
\eu^x\cos y = \Bigl(1 + x + \frac{x^2}{2} +
o(x^2)\Bigr)\Bigl(1 - \frac{y^2}{2} + o(y^2)\Bigr)
= 1 + x + \frac{x^2 - y^2}{2} + o\bigl(\norm{(x,y)}^2\bigr) .
\]
\emph{Por derivadas parciais:} $f_x = \eu^x\cos y$, $f_y =
-\eu^x\sin y$, logo $\nabla f(0) = (1, 0)$; e $f_{xx} = f$,
$f_{yy} = -f$, $f_{xy} = -\eu^x\sin y$ dão $H_0 =
\operatorname{diag}(1, -1)$: \cref{thm:b2:diffcalc:taylor}
reproduz $1 + x + \frac12(x^2 - y^2)$. Os dois cálculos
coincidem, e a rota por composição foi mais rápida --- nenhuma parcial
segunda. Lição final: a origem \emph{não} é um
ponto crítico ($\nabla f \neq 0$), de modo que, apesar da
hessiana indefinida, não há sela a declarar: o termo linear manda,
e o teste de segunda ordem só fala em pontos
críticos.
\end{example}

\begin{theorem}[Teste de extremo de segunda ordem, demonstrado]\label{thm:b2:diffcalc:extrema}
Seja $f$ de classe $C^2$ perto de um ponto crítico $a$ ($\nabla f(a) = 0$),
com hessiana $H = H_a$.
\begin{enumerate}
  \item Se $H$ é positiva definida, $a$ é mínimo local estrito
        (negativa definida: máximo).
  \item Se $H$ tem \omterm{def:b2:reduction:eigen}{autovalores} dos dois sinais, $a$ é uma sela: sem
        extremo.
  \item Se $H$ é singular (e semidefinida), nada se conclui.
\end{enumerate}
O teste ``$rt - s^2$'' do primeiro ano é o caso $n = 2$: $\det H = rt -
s^2$, com o sinal de $\operatorname{tr}$ lido em $r$.
\end{theorem}

\begin{proof}
Pelo teorema espectral (\cref{thm:b2:quadratic:spectral}), a \omterm{def:b2:quadratic:def}{forma quadrática}
de $H$ fica espremida entre seus \omterm{def:b2:reduction:eigen}{autovalores} extremos:
$\lambda_{\min}\norm h^2 \leq \langle Hh, h\rangle \leq
\lambda_{\max}\norm h^2$.

(1) Se $\lambda_{\min} > 0$: Taylor--Young dá
\[
f(a + h) - f(a) \geq \frac{\lambda_{\min}}{2}\norm h^2 -
o(\norm h^2) > 0
\]
para $h \neq 0$ pequeno: mínimo local estrito.

(2) Ao longo de um \omterm{def:b2:reduction:eigen}{autovetor} $v_+$ com $\lambda_+ > 0$: $f(a + tv_+) -
f(a) = \frac{\lambda_+}{2}t^2 + o(t^2) > 0$ para $t$ pequeno; ao longo de
$v_-$ com $\lambda_- < 0$ a diferença é negativa: os dois sinais
ocorrem em toda vizinhança.

(3) $f(x,y) = x^2 + y^4$, $x^2 - y^4$, $x^2 + y^3$ compartilham a mesma
hessiana singular semidefinida em $0$ com três comportamentos
diferentes.
\end{proof}

\begin{example}[Uma classificação completa, global inclusive]\label{ex:b2:diffcalc:quarticrun}
Classifique todos os extremos de $f(x, y) = x^4 + y^4 - 4xy$ em $\R^2$.
\emph{Pontos críticos:} $\nabla f = (4x^3 - 4y,\ 4y^3 - 4x) =
0$ dá $y = x^3$ e $x = y^3 = x^9$, logo $x(x^8 - 1) = 0$:
as soluções reais são $(0,0)$, $(1,1)$, $(-1,-1)$.
\emph{Hessianas:} $H = \begin{pmatrix} 12x^2 & -4\\ -4 &
12y^2\end{pmatrix}$. Em $(\pm1, \pm1)$: $\begin{pmatrix} 12 &
-4\\ -4 & 12\end{pmatrix}$, \omterm{def:b2:reduction:eigen}{autovalores} $8$ e $16$: positiva
definida, mínimos locais estritos com $f = -2$. Em $(0,0)$:
$\begin{pmatrix} 0 & -4\\ -4 & 0\end{pmatrix}$, \omterm{def:b2:reduction:eigen}{autovalores}
$\pm4$: uma sela. \emph{Globalidade:} de $2\abs{xy} \leq x^2 +
y^2$,
\[
f(x, y) \geq x^4 + y^4 - 2(x^2 + y^2)
= (x^2 - 1)^2 + (y^2 - 1)^2 + x^2 + y^2 - 2
\xrightarrow[\norm{(x,y)}\to\infty]{} +\infty :
\]
$f$ é coerciva, logo atinge um mínimo global
(\omterm{def:b2:metric:compact}{compacidade} dos conjuntos de subnível), necessariamente num ponto
crítico: o valor $-2$, tanto em $(1,1)$ quanto em $(-1,-1)$, é o
mínimo global; não há máximo ($f$ é ilimitada superiormente).
Lição final: o teste local classifica candidatos, mas
só um argumento de crescimento transforma ``local'' em ``global'' --- o
padrão em dois passos de toda demonstração de otimização deste livro.
\end{example}

\begin{method}[Classificar os extremos de $f \colon \R^n \to \R$]\label{met:b2:diffcalc:extrema}
\begin{enumerate}
  \item Resolva $\nabla f = 0$ (todos os pontos críticos; num domínio
        com fronteira, trate a fronteira à parte, como no
        \cref{exo:b2:diffcalc:7}).
  \item Em cada ponto crítico, calcule a hessiana e os sinais de seus
        \omterm{def:b2:reduction:eigen}{autovalores} --- em dimensão $2$, apenas $\det H$
        e $\operatorname{tr} H$: $\det < 0$ sela; $\det >
        0$ extremo, do tipo dado pelo sinal do
        traço; $\det = 0$: o teste se cala, estude $f$
        ao longo de curvas.
  \item Para enunciados globais, acrescente um argumento de \omterm{def:b2:metric:compact}{compacidade} ou de coercividade
        ($f \to +\infty$ no infinito, ou um conjunto de restrições
        \omterm{def:b2:metric:compact}{compacto}) e depois compare os valores críticos.
\end{enumerate}
\end{method}

\section{O teorema da função inversa}

\begin{theorem}[Teorema da função inversa]\label{thm:b2:diffcalc:inverse}
Seja $f \colon U \to \R^n$ de classe $C^1$ e $a \in U$ com $\dd f_a$
\emph{invertível}. Então existem vizinhanças \omterm{def:b2:metric:topology}{abertas} $V \ni a$, $W
\ni f(a)$ tais que $f \colon V \to W$ é uma bijeção com inversa
$C^1$, e\index{teorema da função inversa}
\[
\dd (f^{-1})_{f(x)} = (\dd f_x)^{-1} \qquad (x \in V).
\]
\end{theorem}
\admitted

\begin{remark}[Por que é verdade: a estratégia do ponto fixo]
Resolver $f(x) = y$ perto de $a$ reescreve-se como a equação de ponto fixo
$x = x + \dd f_a^{-1}\bigl(y - f(x)\bigr) =: \Phi_y(x)$; a aplicação
$\Phi_y$ tem \omterm{def:b2:diffcalc:differential}{diferencial} $\mathrm{id} - \dd f_a^{-1}\dd f_x$,
pequena perto de $a$ pela \omterm{def:b2:metric:continuity}{continuidade} de $\dd f$, de modo que $\Phi_y$ é uma
contração numa bola fechada pequena e o teorema do ponto fixo de Banach
(\cref{thm:b2:metric:banach}) fornece a única solução local
$x = f^{-1}(y)$. A \omterm{def:b2:metric:continuity}{continuidade} e a \omterm{def:b2:diffcalc:differential}{diferenciabilidade} da
inversa decorrem então das estimativas da contração. A contabilidade \omterm{def:b2:metric:complete}{completa}
é executada no terceiro ano; a estratégia --- e o
enunciado --- são usados livremente daqui em diante. O
\emph{teorema da função implícita} companheiro (resolver $F(x, y) = 0$ em $y(x)$
quando $\frac{\partial F}{\partial y}$ é invertível) segue
aplicando o teorema a $(x, y) \mapsto (x, F(x,y))$.
\end{remark}

\begin{example}[Coordenadas polares]\label{ex:b2:diffcalc:polar}
$\Phi(r, \theta) = (r\cos\theta, r\sin\theta)$ tem jacobiana
\[
J_\Phi = \begin{pmatrix} \cos\theta & -r\sin\theta\\ \sin\theta &
r\cos\theta \end{pmatrix},
\qquad \det J_\Phi = r :
\]
invertível para $r \neq 0$, de modo que $\Phi$ é um difeomorfismo local
$C^1$ fora da origem --- a licença para ``passar a coordenadas
polares'', renovada para as integrais múltiplas do
\cref{ch:b2:multint}.
\end{example}

\begin{example}[Local em toda parte, global em lugar nenhum]\label{ex:b2:diffcalc:localnotglobal}
Seja $f(x, y) = \bigl(\eu^x\cos y,\ \eu^x\sin y\bigr)$ em
$\R^2$. Sua jacobiana,
\[
J_f = \begin{pmatrix}
\eu^x\cos y & -\eu^x\sin y\\
\eu^x\sin y & \eu^x\cos y
\end{pmatrix},
\qquad
\det J_f = \eu^{2x} > 0 ,
\]
nunca se anula: pelo~\cref{thm:b2:diffcalc:inverse}, $f$ é um
difeomorfismo local $C^1$ em \emph{todo} ponto do plano.
E no entanto $f$ está longe de ser injetiva: $f(x, y + 2\pi) = f(x, y)$, de modo que
todo valor é assumido uma infinidade de vezes; e ela também não é sobrejetiva,
pois $\norm{f(x, y)} = \eu^x > 0$ nunca atinge a origem.
Lição final: o teorema da função inversa é irredutivelmente
\emph{local} --- a invertibilidade de todo $\dd f_a$ produz uma
colcha de retalhos de inversas locais que não precisam se juntar numa só.
(Quem conhece números complexos reconhecerá $z \mapsto
\eu^z$; a colcha é a família de ramos do logaritmo.)
Compare com o~\cref{exo:b2:diffcalc:12}, em que uma hipótese
global quantitativa força de fato uma inversa global.
\end{example}

\begin{remark}[Armadilhas comuns]
\emph{(i) Derivadas direcionais são baratas, \omterm{def:b2:diffcalc:differential}{diferenciais} não
são:} todas as derivadas direcionais podem existir --- e até deixar
de depender linearmente da direção --- sem
\omterm{def:b2:diffcalc:differential}{diferenciabilidade} (\cref{exo:b2:diffcalc:2}); só o critério $C^1$
(\cref{thm:b2:diffcalc:c1}) eleva as parciais a uma
\omterm{def:b2:diffcalc:differential}{diferencial}. \emph{(ii) Crítico não significa extremo:}
as selas (\cref{ex:b2:diffcalc:quarticrun}) e o caso singular
silencioso (\cref{thm:b2:diffcalc:extrema}~(3)) se escondem
os dois atrás de $\nabla f = 0$. \emph{(iii) Não há igualdade do
valor médio para valores vetoriais:} só a \emph{desigualdade} do
\cref{thm:b2:diffcalc:mvi} sobrevive
(\cref{exo:b2:diffcalc:9}); nunca escreva $f(b) - f(a) = \dd
f_c(b-a)$ para $f$ com valores em $\R^m$, $m \geq 2$.
\emph{(iv) A invertibilidade local não é injetividade:}
\cref{ex:b2:diffcalc:localnotglobal}. \emph{(v) O gradiente
pertence ao produto interno:} $\nabla f$ é o vetor
que representa $\dd f_a$ num produto interno escolhido; mude o
produto (como no exemplo ponderado do capítulo de \omterm{def:b2:quadratic:def}{formas
quadráticas}) e o gradiente gira, enquanto a \omterm{def:b2:diffcalc:differential}{diferencial} ---
o objeto intrínseco --- não se move.
\end{remark}

\begin{remark}[Onde isso é usado]
Tudo o que vem depois deste capítulo é cálculo diferencial
aplicado: o capítulo de equações diferenciais lineariza fluxos e
usa a fórmula do \omterm{def:b2:linalg:det}{determinante} de Liouville (demonstrada no problema
de fim de semana deste capítulo); os capítulos sobre curvas e superfícies estudam
conjuntos de nível e parametrizações pelo teorema da função
implícita; as integrais múltiplas mudam de variáveis por jacobianos.
O problema de fim de semana desenvolve o cálculo \emph{no próprio espaço
das matrizes} --- \omterm{def:b2:diffcalc:differential}{diferencial} do \omterm{def:b2:linalg:det}{determinante}, da
inversa, a \omterm{ex:b2:nvs:matrixexp}{exponencial de matriz} e o grupo ortogonal como conjunto de nível
suave --- a sombra, no segundo ano, do que o volume do terceiro ano
formaliza como variedades e grupos de Lie.
\end{remark}

\section{Exercícios}

\begin{exercise}[$\star$]\label{exo:b2:diffcalc:1}
Calcule as \omterm{def:b2:diffcalc:differential}{matrizes jacobianas} de $f(x,y) = (x^2 - y^2,\, 2xy)$ e
de $\Phi(r,\theta,z) = (r\cos\theta, r\sin\theta, z)$; onde são
invertíveis as \omterm{def:b2:diffcalc:differential}{diferenciais}?
\end{exercise}

\begin{exercise}[$\star$]\label{exo:b2:diffcalc:2}
Seja $f(x,y) = \frac{x^3}{x^2 + y^2}$ ($f(0,0) = 0$). Prove que todas as
derivadas direcionais de $f$ em $0$ existem, mas que $f$ não é
\omterm{def:b2:diffcalc:differential}{diferenciável} em $0$ \emph{(a aplicação $v \mapsto$ derivada
direcional não é linear)}.
\end{exercise}

\begin{exercise}[$\star$]\label{exo:b2:diffcalc:3}
Ache e classifique os pontos críticos de $f(x, y) = x^3 + y^3 -
3xy$ usando \cref{thm:b2:diffcalc:extrema}, e de $g(x,y) = x^4 +
y^4 - 2(x - y)^2$.
\end{exercise}

\begin{exercise}[$\star\star$]\label{exo:b2:diffcalc:4}
Seja $f \colon \R^n \to \R$ de classe $C^1$ e \emph{homogênea de grau
$p$}: $f(tx) = t^pf(x)$ para $t > 0$. Demonstre a identidade de Euler
\[
\langle \nabla f(x), x\rangle = p\,f(x) ,
\]
e sua recíproca para funções $C^1$ em $\R^n\setminus\{0\}$.
\end{exercise}

\begin{exercise}[$\star\star$]\label{exo:b2:diffcalc:5}
Sejam $A$ \omterm{def:b2:quadratic:adjoint}{simétrica} e $f(x) = \frac12\langle Ax, x\rangle -
\langle b, x\rangle$. Calcule $\nabla f$ e $H_f$; quando $f$ é
convexa? Supondo $A$ positiva definida, mostre que $f$ tem um único
mínimo global na solução de $Ax = b$ --- a razão de ser do
método do gradiente.
\end{exercise}

\begin{exercise}[$\star\star$]\label{exo:b2:diffcalc:6}
(Multiplicador de Lagrange, uma restrição, demonstrado à mão) Sejam $f, g$
de classe $C^1$ em $\R^2$ e suponha que $f$ atinja, em $a$, um extremo local
no conjunto de nível $\{g = 0\}$, com $\nabla g(a) \neq 0$. Prove que
$\nabla f(a) = \lambda\nabla g(a)$ para algum $\lambda$.
\emph{(Parametrize o conjunto de nível perto de $a$ pelo teorema da função
implícita e derive $t \mapsto f(\gamma(t))$.)} Aplicação:
extremos de $f(x,y) = xy$ no círculo $x^2 + y^2 = 1$.
\end{exercise}

\begin{exercise}[$\star\star$]\label{exo:b2:diffcalc:7}
Determine os extremos de $f(x, y) = x^2 + y^2 - xy + x - y$ em
$\R^2$, e depois seu máximo e mínimo no triângulo fechado de
vértices $(0,0)$, $(1,0)$, $(0,1)$ \emph{(pontos críticos interiores,
depois os três lados, depois os vértices)}.
\end{exercise}

\begin{exercise}[$\star\star\star$]\label{exo:b2:diffcalc:8}
Sejam $f \colon \R^2 \to \R^2$, $f(x, y) = (x + y^2,\; y + x^2)$.
Mostre que $f$ é um difeomorfismo local perto de $0$, calcule
$\dd(f^{-1})_{(0,0)}$ e ache o maior $r$ tal que $\dd f$
seja invertível na bola $\norm{(x,y)}_2 < r$ \emph{(calcule
$\det J_f$)}.
\end{exercise}

\begin{exercise}[$\star\star\star$]\label{exo:b2:diffcalc:9}
(Rolle falha, o valor médio sobrevive) Dê $f \colon \R \to \R^2$,
$C^1$, com $f(0) = f(2\pi)$ mas $f'(t) \neq 0$ para todo $t$
(sem Rolle com valores vetoriais). Depois verifique em seu exemplo a
\emph{desigualdade} do valor médio do~\cref{thm:b2:diffcalc:mvi}.
\end{exercise}

\begin{exercise}[$\star$]\label{exo:b2:diffcalc:10}
Calcule a \omterm{def:b2:diffcalc:differential}{diferencial} e o gradiente de $f(x) = \norm
x_2^2$ e de $g(x) = \langle Ax, x\rangle$ em $\R^n$ ($A$ uma
matriz quadrada, não suposta \omterm{def:b2:quadratic:adjoint}{simétrica}) e a hessiana de cada uma.
Para quais $A$ a função $g$ é convexa?
\end{exercise}

\begin{exercise}[$\star\star$]\label{exo:b2:diffcalc:11}
Seja $f \colon \R^n \to \R$ de classe $C^2$. Prove que $f$ é convexa se
e somente se sua hessiana $H_x$ é positiva semidefinida em todo
$x$ \emph{(reduza a uma variável: $t \mapsto f(a + t(b-a))$;
use Taylor--Lagrange numa direção e, para a recíproca,
avalie $\varphi''$)}.
\end{exercise}

\begin{exercise}[$\star\star\star$]\label{exo:b2:diffcalc:12}
(Um teorema de inversão global) Seja $g \colon \R^n \to \R^n$ de classe $C^1$
com $\vertiii{\dd g_x} \leq k < 1$ para todo $x$, e $f =
\mathrm{id} + g$.
\begin{enumerate}
  \item Mostre que $\norm{f(x) - f(y)} \geq (1 - k)\norm{x - y}$: $f$
        é injetiva, com inversa \omterm{def:b2:metric:continuity}{contínua} em sua imagem.
  \item Mostre que, para cada $y \in \R^n$, a aplicação $x \mapsto y -
        g(x)$ é uma contração do espaço \omterm{def:b2:metric:complete}{completo} $\R^n$,
        e conclua pelo teorema do ponto fixo de Banach
        (\cref{thm:b2:metric:banach}) que $f$ é sobrejetiva.
  \item Conclua que $f$ é uma bijeção de $\R^n$ com
        inversa $(1-k)^{-1}$-lipschitziana --- uma contrapartida
        \emph{global} do~\cref{thm:b2:diffcalc:inverse} (que,
        em contraste, é puramente local).
\end{enumerate}
\end{exercise}

\section{Problema: o cálculo das matrizes --- Jacobi,
exponencial e o grupo ortogonal}

\begin{problem}\label{pb:b2:diffcalc:1}
O playground mais limpo para o cálculo diferencial é o próprio espaço
$\mathcal M_n(\R) \simeq \R^{n^2}$: suas aplicações mais naturais
--- produto, inversa, \omterm{def:b2:linalg:det}{determinante}, exponencial --- têm
\omterm{def:b2:diffcalc:differential}{diferenciais} de elegância marcante. Este problema as calcula
todas: a \emph{série de Neumann}, a \omterm{def:b2:diffcalc:differential}{diferencial} da inversa,
a \emph{fórmula de Jacobi}\index{fórmula de Jacobi} para o
\omterm{def:b2:linalg:det}{determinante} com a \emph{fórmula de Liouville} como dividendo, a
\emph{exponencial de matriz}\index{exponencial de matriz} com
$\det\eu^A = \eu^{\operatorname{tr}A}$, e por fim o
grupo ortogonal $O_n$ como conjunto de nível suave com as
matrizes antissimétricas como espaço tangente --- geometria
diferencial em embrião. Em todo o problema, $\vertiii\cdot$ é a \omterm{thm:b2:nvs:continuouslinear}{norma de
operador} subordinada a $\norm\cdot_2$, e $\langle X, Y\rangle =
\operatorname{tr}(X^{\mathsf T}Y)$ o produto interno de Frobenius.

\medskip
\noindent\textbf{Parte I --- A série de Neumann.}
\begin{enumerate}
  \item Demonstre a submultiplicatividade, $\vertiii{AB} \leq
        \vertiii A\,\vertiii B$, e deduza que as aplicações
        polinomiais de $A$ (produtos de matrizes, \omterm{def:b2:linalg:det}{determinante}, traço) são
        \omterm{def:b2:metric:continuity}{contínuas} em $\mathcal M_n(\R)$.
  \item Para $\vertiii X < 1$, mostre que $\sum_{k\geq0}X^k$
        converge \omterm{def:b2:series:def}{absolutamente} em $\mathcal M_n(\R)$
        (\cref{thm:b2:nvs:absoluteconvergence}), que sua soma
        é $(I - X)^{-1}$ e que
        \[
        \vertiii{(I - X)^{-1}} \leq \frac{1}{1 - \vertiii X},
        \qquad
        (I - X)^{-1} = I + X + O\bigl(\vertiii X^2\bigr) .
        \]
  \item Deduza que $GL_n(\R)$ é \emph{\omterm{def:b2:metric:topology}{aberto}}: se $A$ é
        invertível e $\vertiii H <
        \frac{1}{\vertiii{A^{-1}}}$, então $A + H$ é
        invertível. Deduza também que $GL_n(\R)$ é \emph{denso}
        em $\mathcal M_n(\R)$ \emph{(perturbe $A$ por
        $\varepsilon I$: $\det(A + \varepsilon I)$ é um polinômio não nulo
        em $\varepsilon$)}.
  \item Mostre que a aplicação inversão $\Phi(A) = A^{-1}$ é
        \omterm{def:b2:metric:continuity}{contínua} em $GL_n(\R)$.
\end{enumerate}

\medskip
\noindent\textbf{Parte II --- Primeiras \omterm{def:b2:diffcalc:differential}{diferenciais}.}
\begin{enumerate}[resume]
  \item Mostre que a aplicação quadrado $A \mapsto A^2$ é
        \omterm{def:b2:diffcalc:differential}{diferenciável} com \omterm{def:b2:diffcalc:differential}{diferencial} $H \mapsto AH + HA$,
        e, mais geralmente, que $A \mapsto A^k$ tem
        \omterm{def:b2:diffcalc:differential}{diferencial} $H \mapsto \sum_{i=0}^{k-1}
        A^iHA^{k-1-i}$. Por que \emph{não} se pode escrever $kA^{k-1}H$
        em geral?
  \item Prove que $\Phi(A) = A^{-1}$ é \omterm{def:b2:diffcalc:differential}{diferenciável} em
        $GL_n(\R)$ com
        \[
        \dd\Phi_A(H) = -A^{-1}HA^{-1}
        \]
        \emph{(escreva $(A + H)^{-1} = (I + A^{-1}H)^{-1}A^{-1}$
        e desenvolva pela questão 2)}. Confira a fórmula contra
        o caso escalar $n = 1$.
  \item Para uma curva $t \mapsto A(t) \in GL_n(\R)$ de classe $C^1$, deduza que
        $\bigl(A(t)^{-1}\bigr)' = -A^{-1}A'A^{-1}$ e desenvolva
        $t \mapsto (I + tB)^{-1}$ até a primeira ordem em $t = 0$.
  \item Calcule a \omterm{def:b2:diffcalc:differential}{diferencial} de $f(A) =
        \operatorname{tr}(A^k)$ e identifique seu gradiente para
        o produto interno de Frobenius:
        \[
        \dd f_A(H) = k\operatorname{tr}\bigl(A^{k-1}H\bigr),
        \qquad
        \nabla f(A) = k\,\bigl(A^{k-1}\bigr)^{\mathsf T} .
        \]
  \item Mesmas questões para $f(A) = \operatorname{tr}
        (A^{\mathsf T}A) = \norm A_F^2$: \omterm{def:b2:diffcalc:differential}{diferencial},
        gradiente e a hessiana (constante); conclua que
        $\norm\cdot_F^2$ é estritamente convexa.
\end{enumerate}

\medskip
\noindent\textbf{Parte III --- A fórmula de Jacobi.}
\begin{enumerate}[resume]
  \item Prove que
        \[
        \det(I + H) = 1 + \operatorname{tr}H +
        O\bigl(\vertiii H^2\bigr)
        \]
        \emph{(desenvolva $\det(e_1 + h_1, \dots, e_n + h_n)$ pela
        multilinearidade nas colunas: os termos com ao menos duas
        colunas de $h$ são $O(\vertiii H^2)$)}: $\dd(\det)_I =
        \operatorname{tr}$.
  \item Para $A$ invertível, deduza que
        \[
        \dd(\det)_A(H) = \det(A)\,
        \operatorname{tr}\bigl(A^{-1}H\bigr) .
        \]
  \item Mostre que, para \emph{toda} $A$ (invertível ou não),
        $\frac{\partial\det}{\partial a_{ij}}(A) = C_{ij}$, o
        cofator $(i,j)$ \emph{(desenvolvimento de Laplace ao longo da linha
        $i$)}, de modo que, com a adjunta $\operatorname{adj}A =
        \operatorname{com}(A)^{\mathsf T}$:
        \[
        \dd(\det)_A(H) =
        \operatorname{tr}\bigl(\operatorname{adj}(A)\,H\bigr),
        \qquad
        \nabla(\det)(A) = \operatorname{com}(A) ,
        \]
        recuperando a questão 11 quando $A$ é invertível
        ($\operatorname{adj}A = \det(A)A^{-1}$). Essa é a
        \emph{fórmula de Jacobi}: $\bigl(\det
        A(t)\bigr)' = \operatorname{tr}\bigl(
        \operatorname{adj}(A(t))\,A'(t)\bigr)$.
  \item (Fórmula de Liouville) Seja $A(t)$ uma curva $C^1$ de
        matrizes satisfazendo a equação diferencial linear
        $A'(t) = M(t)A(t)$. Prove que
        \[
        \bigl(\det A(t)\bigr)' =
        \operatorname{tr}\bigl(M(t)\bigr)\,\det A(t),
        \qquad\text{logo}\qquad
        \det A(t) = \det A(0)\,
        \exp\Bigl(\int_0^t\operatorname{tr}M\Bigr)
        \]
        \emph{(use $\operatorname{adj}(A)\,A = \det(A)I$ e a
        invariância \omterm{def:b2:structures:generated}{cíclica} do traço)} --- a identidade do wronskiano
        que o capítulo de equações diferenciais usará
        constantemente.
  \item Mostre que $SL_n(\R) = \{\det = 1\}$ é um conjunto de nível
        suave: em todo $A \in SL_n(\R)$ a \omterm{def:b2:diffcalc:differential}{diferencial}
        $\dd(\det)_A$ é uma aplicação linear \emph{sobrejetiva} sobre
        $\R$ \emph{(avalie-a em $H = \frac1nA$)}.
\end{enumerate}

\medskip
\noindent\textbf{Parte IV --- A \omterm{ex:b2:nvs:matrixexp}{exponencial de matriz}.}
\begin{enumerate}[resume]
  \item Mostre que $\eu^A = \sum_{k\geq0}\frac{A^k}{k!}$
        converge \omterm{def:b2:series:def}{absolutamente} para toda $A$, \omterm{def:b2:funcseq:series}{normalmente} em toda
        bola, com $\vertiii{\eu^A} \leq
        \eu^{\vertiii A}$; e que $\eu^A$ depende
        \omterm{def:b2:metric:continuity}{continuamente} de $A$.
  \item Prove que $AB = BA$ implica $\eu^{A+B} =
        \eu^A\eu^B$ \emph{(\omterm{thm:b2:series:fubini}{produto de Cauchy}, legítimo pela
        convergência absoluta)}; deduza que $\eu^A$ é sempre
        invertível, de inversa $\eu^{-A}$: $\exp$ leva
        $\mathcal M_n(\R)$ em $GL_n(\R)$.
  \item Mostre que $t \mapsto \eu^{tA}$ é $C^1$ (na verdade
        $C^\infty$) com
        \[
        \frac{\dd}{\dd t}\,\eu^{tA} = A\,\eu^{tA} =
        \eu^{tA}A
        \]
        \emph{(derive a série termo a termo em segmentos: a
        série derivada converge \omterm{def:b2:funcseq:series}{normalmente})}.
  \item Demonstre a identidade
        \[
        \det\bigl(\eu^{A}\bigr) = \eu^{\operatorname{tr}A}
        \]
        \emph{(aplique a fórmula de Liouville, questão 13, a
        $A(t) = \eu^{tA}$)}. Verificações de bom senso: $n = 1$; $A$
        nilpotente; e as matrizes de traço nulo caem em
        $SL_n(\R)$.
  \item Mostre que $\eu^H = I + H + O(\vertiii H^2)$, de modo que $\exp$ é
        \omterm{def:b2:diffcalc:differential}{diferenciável} em $0$ com $\dd(\exp)_0 =
        \mathrm{id}$; conclua pelo teorema da função
        inversa (\cref{thm:b2:diffcalc:inverse}) que $\exp$
        é um difeomorfismo $C^1$ de uma vizinhança de $0$
        sobre uma vizinhança de $I$: toda matriz próxima da
        identidade tem um logaritmo.
  \item Mostre que $\exp$ leva as matrizes \omterm{def:b2:quadratic:adjoint}{simétricas} nas \omterm{def:b2:quadratic:adjoint}{simétricas}
        \emph{positivas definidas}, bijetivamente
        \emph{(diagonalize; a inversa é o logaritmo
        espectral)}.
\end{enumerate}

\medskip
\noindent\textbf{Parte V --- O grupo ortogonal como conjunto de
nível.}
\begin{enumerate}[resume]
  \item Seja $F(A) = A^{\mathsf T}A$, de $\mathcal M_n(\R)$ nas
        matrizes \omterm{def:b2:quadratic:adjoint}{simétricas} $S_n$. Calcule $\dd F_A(H) =
        A^{\mathsf T}H + H^{\mathsf T}A$ e mostre que, em todo
        $A \in O_n = F^{-1}(I)$, essa \omterm{def:b2:diffcalc:differential}{diferencial} é
        \emph{sobrejetiva} sobre $S_n$ \emph{(dada $S \in S_n$,
        tente $H = \frac12 AS$)}: $O_n$ é um conjunto de nível suave,
        de dimensão $n^2 - \frac{n(n+1)}2 = \frac{n(n-1)}2$.
  \item Mostre que toda curva $A(t) \in O_n$ de classe $C^1$ com $A(0) =
        I$ tem velocidade \emph{antissimétrica} $A'(0)$ e,
        reciprocamente, que para $K$ antissimétrica a curva
        $\eu^{tK}$ permanece em $O_n$: o espaço tangente de $O_n$
        em $I$ é exatamente o das matrizes antissimétricas.
  \item Mostre que $\det\eu^{K} = 1$ para $K$ antissimétrica (questão
        18): a curva exponencial vive no grupo de rotações
        $SO_n$. Calcule-a por inteiro para $n = 2$: com $J =
        \begin{pmatrix}0 & -1\\ 1 & 0\end{pmatrix}$, prove que
        \[
        \eu^{\theta J} = \begin{pmatrix} \cos\theta &
        -\sin\theta\\ \sin\theta & \cos\theta\end{pmatrix} :
        \]
        a \omterm{ex:b2:nvs:matrixexp}{exponencial de matriz} \emph{é} a rotação de
        $\theta$, e as definições em série do cosseno e do seno
        reaparecem dentro de uma matriz.
  \item (O truque da densidade) Usando a densidade de $GL_n(\R)$
        (questão 3) e a \omterm{def:b2:metric:continuity}{continuidade}, estenda das invertíveis a
        todas as matrizes a identidade
        \[
        \operatorname{adj}(AB) =
        \operatorname{adj}(B)\operatorname{adj}(A)
        \]
        \emph{(para $A, B$ invertível os dois lados valem
        $\det(AB)(AB)^{-1}$; os dois lados são polinomiais nas
        entradas)}.
  \item Síntese. Uma frase para cada: (i) quais capítulos anteriores
        forneceram o motor de cada Parte
        (a \omterm{def:b2:metric:complete}{completude} e as \omterm{def:b2:structures:algebra}{álgebras} normadas; o teorema
        espectral; o teorema da função inversa); (ii) qual
        fórmula deste problema o capítulo de equações
        diferenciais vai usar, e onde; (iii) o que
        $\dd(\det)_I = \operatorname{tr}$ e $\det\eu^A =
        \eu^{\operatorname{tr}A}$ dizem sobre o traço e o
        \omterm{def:b2:linalg:det}{determinante} como ``volume infinitesimal e global'';
        (iv) o que o volume do terceiro ano de graduação faz das questões
        21--23 (grupos de Lie e suas \omterm{def:b2:structures:algebra}{álgebras} de Lie).
\end{enumerate}
\end{problem}
